% Options for packages loaded elsewhere
\PassOptionsToPackage{unicode}{hyperref}
\PassOptionsToPackage{hyphens}{url}
%
\documentclass[
  10pt,
]{article}
\usepackage{amsmath,amssymb}
\usepackage{iftex}
\ifPDFTeX
  \usepackage[T1]{fontenc}
  \usepackage[utf8]{inputenc}
  \usepackage{textcomp} % provide euro and other symbols
\else % if luatex or xetex
  \usepackage{unicode-math} % this also loads fontspec
  \defaultfontfeatures{Scale=MatchLowercase}
  \defaultfontfeatures[\rmfamily]{Ligatures=TeX,Scale=1}
\fi
\usepackage{lmodern}
\ifPDFTeX\else
  % xetex/luatex font selection
\fi
% Use upquote if available, for straight quotes in verbatim environments
\IfFileExists{upquote.sty}{\usepackage{upquote}}{}
\IfFileExists{microtype.sty}{% use microtype if available
  \usepackage[]{microtype}
  \UseMicrotypeSet[protrusion]{basicmath} % disable protrusion for tt fonts
}{}
\makeatletter
\@ifundefined{KOMAClassName}{% if non-KOMA class
  \IfFileExists{parskip.sty}{%
    \usepackage{parskip}
  }{% else
    \setlength{\parindent}{0pt}
    \setlength{\parskip}{6pt plus 2pt minus 1pt}}
}{% if KOMA class
  \KOMAoptions{parskip=half}}
\makeatother
\usepackage{xcolor}
\usepackage[margin=1in]{geometry}
\usepackage{listings}
\newcommand{\passthrough}[1]{#1}
\lstset{defaultdialect=[5.3]Lua}
\lstset{defaultdialect=[x86masm]Assembler}
\usepackage{longtable,booktabs,array}
\usepackage{calc} % for calculating minipage widths
% Correct order of tables after \paragraph or \subparagraph
\usepackage{etoolbox}
\makeatletter
\patchcmd\longtable{\par}{\if@noskipsec\mbox{}\fi\par}{}{}
\makeatother
% Allow footnotes in longtable head/foot
\IfFileExists{footnotehyper.sty}{\usepackage{footnotehyper}}{\usepackage{footnote}}
\makesavenoteenv{longtable}
\setlength{\emergencystretch}{3em} % prevent overfull lines
\providecommand{\tightlist}{%
  \setlength{\itemsep}{0pt}\setlength{\parskip}{0pt}}
\setcounter{secnumdepth}{5}
\ifLuaTeX
  \usepackage{selnolig}  % disable illegal ligatures
\fi
\usepackage{bookmark}
\IfFileExists{xurl.sty}{\usepackage{xurl}}{} % add URL line breaks if available
\urlstyle{same}
\hypersetup{
  pdftitle={XSTAR Atomic Database, Source Architecture, Physics, and Implementation Roadmap},
  pdfauthor={xstar-atomic development documentation},
  hidelinks,
  pdfcreator={LaTeX via pandoc}}

\title{XSTAR Atomic Database, Source Architecture, Physics, and
Implementation Roadmap}
\author{xstar-atomic development documentation}
\date{Version 0.3.200}

\begin{document}
\maketitle

{
\setcounter{tocdepth}{3}
\tableofcontents
}
\section{XSTAR atomic database, source architecture, physics, and
implementation
roadmap}\label{xstar-atomic-database-source-architecture-physics-and-implementation-roadmap}

\textbf{Document version:} 0.3.200\\
\textbf{Package:} \passthrough{\lstinline!xstar-atomic!}\\
\textbf{Primary validation case:} the same-run He-like C V, O VII, Mg
XI, and Ca XIX benchmark suite

\subsection{1. Purpose and scope}\label{purpose-and-scope}

This document is the technical bridge between three representations of
the same calculation:

\begin{enumerate}
\def\labelenumi{\arabic{enumi}.}
\tightlist
\item
  the packed records in XSTAR's \passthrough{\lstinline!atdb.fits!}
  atomic database;
\item
  the XSTAR Fortran routines that interpret those records and assemble
  local rate equations;
\item
  the Python objects, audit products, and prototype solvers in
  \passthrough{\lstinline!xstar-atomic!}.
\end{enumerate}

It is intended both as a reference for understanding XSTAR and as an
implementation plan for a source-code-equivalent Python package with a
later C++ backend. It records the findings through
\passthrough{\lstinline!xstar-atomic!} v0.3.199, including the direct
\passthrough{\lstinline!ucalc!}, matrix-insertion, population-vector,
compact-basis, and row-balance probes.

The current package is not yet a complete replacement for XSTAR. It can
decode the database, evaluate and audit many individual processes,
reconstruct selected local matrices, and compare them with instrumented
XSTAR. A complete implementation still requires the full element basis,
native source and normalization closure, thermal/ionization iteration,
radial transfer, and output writers.

\subsection{2. Executive model of how XSTAR
works}\label{executive-model-of-how-xstar-works}

For each radial zone, XSTAR performs a coupled calculation:

\begin{lstlisting}
input spectrum and shell state
  -> local transferred continuum (epi, bremsa, bremsint, depths)
  -> temperature/electron-fraction iteration
  -> ionization rates and element compact population matrices
  -> level and superlevel populations
  -> line, RRC, continuum emissivity and opacity
  -> heating/cooling and optical-depth updates
  -> next zone/pass
  -> xout_* and xoNN_detail/detal* products
\end{lstlisting}

The central source-code path is:

\begin{lstlisting}
src/xstar/xstar.f90
  -> xstarsetup / readtbl / setptrs
  -> radial shell and pass loop
       -> step
       -> trnfrc
       -> xstarcalc
            -> bremsmap
            -> dsec (when temperature/electron fraction are iterated)
            -> calc_hmc_all
                 -> calc_hmc_element
                      -> levwkelement
                      -> calc_hmc_ion
                           -> ucalc
                      -> msolvelucy
            -> calc_emisab_all
            -> calc_emis_all
       -> heatt / stpcut / savd
  -> pprint / writespectra*
\end{lstlisting}

The important architectural lesson is that final
\passthrough{\lstinline!xout\_*!} products are not a substitute for the
live local arrays used by \passthrough{\lstinline!ucalc!}.
Source-equivalent rates require the same local temperature, density,
populations, optical depths, covering fraction, and radiation grid.

\subsection{\texorpdfstring{3. \texttt{atdb.fits}: physical FITS
layout}{3. atdb.fits: physical FITS layout}}\label{atdb.fits-physical-fits-layout}

\subsubsection{3.1 Primary HDU and four packed vector
extensions}\label{primary-hdu-and-four-packed-vector-extensions}

\passthrough{\lstinline!atdb.fits!} is not organized as one semantic
table per atomic process. The primary HDU stores metadata such as
creation date and creator. Four vector extensions store all records:

\begin{longtable}[]{@{}
  >{\raggedright\arraybackslash}p{(\columnwidth - 4\tabcolsep) * \real{0.3333}}
  >{\raggedright\arraybackslash}p{(\columnwidth - 4\tabcolsep) * \real{0.3333}}
  >{\raggedright\arraybackslash}p{(\columnwidth - 4\tabcolsep) * \real{0.3333}}@{}}
\toprule\noalign{}
\begin{minipage}[b]{\linewidth}\raggedright
FITS extension
\end{minipage} & \begin{minipage}[b]{\linewidth}\raggedright
Fortran/Python representation
\end{minipage} & \begin{minipage}[b]{\linewidth}\raggedright
Contents
\end{minipage} \\
\midrule\noalign{}
\endhead
\bottomrule\noalign{}
\endlastfoot
\passthrough{\lstinline!POINTERS!} &
\passthrough{\lstinline!masterdata\%nptrs!},
\passthrough{\lstinline!ATDB.pointers!} & Ten integers per record
describing type, payload lengths, and offsets. \\
\passthrough{\lstinline!REALS!} &
\passthrough{\lstinline!masterdata\%rdat1!},
\passthrough{\lstinline!ATDB.reals!} & Concatenated single-precision
real payloads. \\
\passthrough{\lstinline!INTEGERS!} &
\passthrough{\lstinline!masterdata\%idat1!},
\passthrough{\lstinline!ATDB.integers!} & Concatenated integer payloads,
including level and ion identifiers. \\
\passthrough{\lstinline!CHARS!} &
\passthrough{\lstinline!masterdata\%kdat1!},
\passthrough{\lstinline!ATDB.chars!} & Concatenated character payloads
such as labels/configurations. \\
\end{longtable}

The FITS \passthrough{\lstinline!LENGTH!} keyword of each extension
gives the logical vector length. XSTAR reads these arrays in
\passthrough{\lstinline!readtbl.f90!};
\passthrough{\lstinline!xstar-atomic!} mirrors this layout in
\passthrough{\lstinline!inspect.py!} and
\passthrough{\lstinline!hierarchy.py!} with memory mapping and lazy
loading of the large real vector.

\subsubsection{3.2 Ten-integer record
header}\label{ten-integer-record-header}

Every record has the following pointer/header block:

\begin{longtable}[]{@{}
  >{\raggedleft\arraybackslash}p{(\columnwidth - 6\tabcolsep) * \real{0.3077}}
  >{\raggedright\arraybackslash}p{(\columnwidth - 6\tabcolsep) * \real{0.2308}}
  >{\raggedright\arraybackslash}p{(\columnwidth - 6\tabcolsep) * \real{0.2308}}
  >{\raggedright\arraybackslash}p{(\columnwidth - 6\tabcolsep) * \real{0.2308}}@{}}
\toprule\noalign{}
\begin{minipage}[b]{\linewidth}\raggedleft
Pointer index
\end{minipage} & \begin{minipage}[b]{\linewidth}\raggedright
Fortran meaning
\end{minipage} & \begin{minipage}[b]{\linewidth}\raggedright
Python field
\end{minipage} & \begin{minipage}[b]{\linewidth}\raggedright
Notes
\end{minipage} \\
\midrule\noalign{}
\endhead
\bottomrule\noalign{}
\endlastfoot
1 & original record pointer/index & \passthrough{\lstinline!raw0!} &
Database-defined identifier. \\
2 & data type (\passthrough{\lstinline!ltyp!}) &
\passthrough{\lstinline!data\_type!} & Selects the payload
schema/formula family. \\
3 & rate type (\passthrough{\lstinline!lrtyp!}) &
\passthrough{\lstinline!rate\_type!} & Selects the physical role and
parent traversal. \\
4 & continuation flag & \passthrough{\lstinline!continuation!} &
Retained for provenance; current \passthrough{\lstinline!dread!} logic
rarely uses it. \\
5 & number of real values & \passthrough{\lstinline!nreal!} & Slice
length in \passthrough{\lstinline!REALS!}. \\
6 & number of integer values & \passthrough{\lstinline!nint!} & Slice
length in \passthrough{\lstinline!INTEGERS!}. \\
7 & number of character bytes & \passthrough{\lstinline!nchar!} & Slice
length in \passthrough{\lstinline!CHARS!}. \\
8 & 1-based pointer into \passthrough{\lstinline!REALS!} &
\passthrough{\lstinline!real\_ptr!} & Python subtracts one before
slicing. \\
9 & 1-based pointer into \passthrough{\lstinline!INTEGERS!} &
\passthrough{\lstinline!int\_ptr!} & Python subtracts one before
slicing. \\
10 & 1-based pointer into \passthrough{\lstinline!CHARS!} &
\passthrough{\lstinline!char\_ptr!} & Python subtracts one before
slicing. \\
\end{longtable}

A raw record is therefore:

\begin{lstlisting}
RecordHeader + real slice + integer slice + character slice
\end{lstlisting}

The record header does not by itself define the semantic meaning of each
payload position. That interpretation is selected by the
\passthrough{\lstinline!(data\_type, rate\_type)!} pair and the relevant
\passthrough{\lstinline!ucalc!}, \passthrough{\lstinline!calt*!}, or
decoder branch.

\subsubsection{3.3 Hierarchical record
order}\label{hierarchical-record-order}

\passthrough{\lstinline!setptrs.f90!} and
\passthrough{\lstinline!xstar\_atomic.hierarchy!} reconstruct the
database hierarchy from record order:

\begin{lstlisting}
element record (rate type 11)
  -> ion records (rate type 12)
       -> level records (rate type 13)
       -> ionization/recombination records
       -> collision records
       -> radiative line records
       -> superlevel and satellite records
\end{lstlisting}

The second-to-last integer in a level or many level-resolved records is
commonly the local level number used by
\passthrough{\lstinline!setptrs!} (\passthrough{\lstinline!nclev!}).
This is a database convention, not a universal promise for every data
type; each decoder must preserve the raw payload and document its
interpretation.

\subsubsection{\texorpdfstring{3.4 Derived pointer structures built by
\texttt{setptrs.f90}}{3.4 Derived pointer structures built by setptrs.f90}}\label{derived-pointer-structures-built-by-setptrs.f90}

XSTAR converts the sequential packed records into navigation arrays. The
most important are:

\begin{longtable}[]{@{}
  >{\raggedright\arraybackslash}p{(\columnwidth - 2\tabcolsep) * \real{0.5000}}
  >{\raggedright\arraybackslash}p{(\columnwidth - 2\tabcolsep) * \real{0.5000}}@{}}
\toprule\noalign{}
\begin{minipage}[b]{\linewidth}\raggedright
Derived pointer
\end{minipage} & \begin{minipage}[b]{\linewidth}\raggedright
Purpose
\end{minipage} \\
\midrule\noalign{}
\endhead
\bottomrule\noalign{}
\endlastfoot
\passthrough{\lstinline!npar(record)!} & Parent element or ion
record. \\
\passthrough{\lstinline!npnxt(record)!} & Next record of the same rate
type under the same parent. \\
\passthrough{\lstinline!npfirst(rate\_type)!} & First database record of
a rate type. \\
\passthrough{\lstinline!npfi(rate\_type, ion)!} & First record of a rate
type for an ion. \\
\passthrough{\lstinline!npfe(rate\_type, element)!} & First record of a
rate type for an element. \\
\passthrough{\lstinline!npilev(level, ion)!} & Database record for a
local level. \\
\passthrough{\lstinline!npilini!} / \passthrough{\lstinline!nplin!}
families & Maps line records and line-emissivity indices. \\
\passthrough{\lstinline!npcon!}, \passthrough{\lstinline!npconi!},
\passthrough{\lstinline!npconi2!} & Connect photoionization records to
continuum-emissivity indices. \\
\passthrough{\lstinline!nlevs(ion)!} & Number of local rows associated
with an ion record, including its compact parent/continuum
convention. \\
\passthrough{\lstinline!npnxt!} chains & Allow
\passthrough{\lstinline!calc\_hmc\_ion!} to iterate records without
scanning the full database. \\
\end{longtable}

\passthrough{\lstinline!xstar-atomic!} preserves the same hierarchy in
\passthrough{\lstinline!ATDBIndexArrays!}: record number, data/rate
types, payload counts and pointers, element, global ion index,
spectroscopic ion stage, local level index, and parent-kind code. The
NPZ hierarchy cache stores numeric arrays so a targeted selection does
not require reconstructing more than a million Python objects.

\subsubsection{3.5 Rate-type catalog}\label{rate-type-catalog}

\begin{longtable}[]{@{}rl@{}}
\toprule\noalign{}
Code & XSTAR rate-type meaning \\
\midrule\noalign{}
\endhead
\bottomrule\noalign{}
\endlastfoot
1 & ground state ionization \\
2 & level ionization/recombination \\
3 & bound-bound collision \\
4 & bound-bound radiative \\
5 & bound-free collision (level) \\
6 & total recombination \\
7 & bound-free radiative (level) \\
8 & total recombination \\
9 & 2 photon decay \\
11 & element data \\
12 & ion data \\
13 & level data \\
14 & radiative superlevel-\textgreater spectroscopic level \\
15 & CI total rate \\
23 & collisional superlevel-\textgreater spectroscopic level \\
40 & CI from superlevels \\
41 & Auger decay / Fe K Auger \\
42 & fluorescence / Auger-related \\
\end{longtable}

\subsubsection{3.6 Data-type catalog}\label{data-type-catalog}

The following labels are those currently carried by
\passthrough{\lstinline!xstar\_atomic.hierarchy!}. Some codes are
historical or database-version dependent; unlabeled codes are
deliberately kept as such instead of assigning an unsupported
interpretation.

\begin{longtable}[]{@{}rl@{}}
\toprule\noalign{}
Code & XSTAR data-type label used by
\passthrough{\lstinline!xstar-atomic!} \\
\midrule\noalign{}
\endhead
\bottomrule\noalign{}
\endlastfoot
1 & radiative recombination: Aldrovandi \& Pequignot \\
2 & charge exchange H0: Kingdon \& Ferland \\
3 & autoionization: Hamilton, Sarazin, Chevalier \\
4 & line data radiative: Mendoza; Raymond \& Smith \\
5 & 2 photon transition collisional \\
6 & level data \\
7 & dielectronic recombination: Aldrovandi \& Pequignot \\
8 & dielectronic recombination: Arnaud \& Raymond \\
9 & charge exchange H0 Kingdon \& Ferland \\
10 & charge exchange H+ Kingdon \& Ferland \\
11 & 2 photon radiative \\
12 & photoionization, excited levels: hydrogenic \\
13 & element data \\
14 & ion data \\
15 & photoionization: Barfield, Koontz \& Huebner \\
16 & Arnaud \& Raymond collisional ionization \\
17 & collisional excitation hydrogenic: Cota \\
18 & radiative recombination hydrogenic: Cota \\
19 & photoionization: HULLAC \\
20 & charge exchange H+ Kingdon \& Ferland \\
21 & PI cross section continued \\
22 & dielectronic recombination: Storey \\
23 & photoionization, excited levels: Clark \\
24 & PI cross section Clark continued \\
25 & collisional ionization: Raymond \& Smith \\
26 & collisional ionization hydrogenic: Cota \\
27 & photoionization: hydrogenic \\
28 & line data collisional: Mendoza; Raymond \& Smith \\
29 & collisional ionization data: scaled hydrogenic \\
30 & radiative recombination hydrogenic: Gould \& Thakur \\
31 & line data no levels \\
32 & collisional ionization: Cota \\
33 & line data collisional: HULLAC \\
34 & line data radiative: Mendoza; Raymond \& Smith \\
35 & photoionization: table from BKH \\
36 & photoionization, excited levels: hydrogenic no level \\
49 & OP PI cross sections for inner shells \\
50 & OP line radiative rates \\
51 & OP and CHIANTI line collisional rates \\
52 & same as 59 but rate type 7 \\
53 & OP PI cross sections \\
54 & H-like Cij, Bautista, H-like ion \\
55 & hydrogenic PI cross sections, Bautista format \\
56 & tabulated collision strength, Bautista \\
57 & effective charge for collisional ionization \\
58 & H-like recombination rates, Bautista \\
59 & Verner PI cross sections \\
60 & Calloway H-like collision strength \\
61 & H-like Cij, Bautista, non-H-like ion \\
62 & Calloway H-like collision strength \\
63 & H-like Cij, Bautista, H-like ion \\
64 & hydrogenic PI cross sections, Bautista format \\
65 & effective charge for collisional ionization \\
66 & like type 69 but fine-structure data \\
67 & effective collision strengths from Keenan et al. \\
68 & He-like collision strengths by Zhang \& Sampson \\
69 & Kato \& Nakazaki fit to He-like collision strengths \\
70 & coefficients for photoionization cross sections of superlevels \\
71 & transition rates from superlevel to spectroscopic levels \\
72 & autoionization rates for satellite levels \\
73 & fit to collisional strengths, satellite levels, He-like ions \\
74 & delta functions added to photoionization cross sections for DR \\
75 & autoionization data for Fe XXIV satellites \\
76 & 2 photon decay \\
77 & collisional rates from 71 \\
78 & Auger level data \\
79 & fluorescence line data \\
80 & collisional ionization rates, ground of Fe and Ni \\
81 & Bhatia Fe XIX collision strengths \\
82 & Fe UTA radiative rates \\
83 & Fe UTA level data \\
84 & Iron K PI cross sections, spectator Auger binned \\
85 & Iron K PI cross sections, spectator Auger summed \\
86 & Iron K Auger data \\
88 & unlabeled XSTAR data type 88 \\
91 & unlabeled XSTAR data type 91 \\
92 & unlabeled XSTAR data type 92 \\
95 & Bryans collisional ionization / CI total rates \\
98 & CHIANTI 2016 collisional excitation rates \\
99 & unlabeled XSTAR data type 99 \\
\end{longtable}

\subsubsection{3.7 Key record layouts used in the current
implementation}\label{key-record-layouts-used-in-the-current-implementation}

\begin{longtable}[]{@{}
  >{\raggedright\arraybackslash}p{(\columnwidth - 6\tabcolsep) * \real{0.2500}}
  >{\raggedright\arraybackslash}p{(\columnwidth - 6\tabcolsep) * \real{0.2500}}
  >{\raggedright\arraybackslash}p{(\columnwidth - 6\tabcolsep) * \real{0.2500}}
  >{\raggedright\arraybackslash}p{(\columnwidth - 6\tabcolsep) * \real{0.2500}}@{}}
\toprule\noalign{}
\begin{minipage}[b]{\linewidth}\raggedright
Data/rate type
\end{minipage} & \begin{minipage}[b]{\linewidth}\raggedright
Process
\end{minipage} & \begin{minipage}[b]{\linewidth}\raggedright
Important payload interpretation
\end{minipage} & \begin{minipage}[b]{\linewidth}\raggedright
Main source routine
\end{minipage} \\
\midrule\noalign{}
\endhead
\bottomrule\noalign{}
\endlastfoot
13/11 & Element metadata & abundance, atomic weight, element identity &
\passthrough{\lstinline!setptrs!},
\passthrough{\lstinline!levwkelement!} \\
14/12 & Ion metadata & global ion index, ion stage, number of levels,
labels & \passthrough{\lstinline!setptrs!},
\passthrough{\lstinline!levwk!},
\passthrough{\lstinline!levwkelement!} \\
6/13, 83/13 & Level data & energy, statistical weight, level number,
configuration & \passthrough{\lstinline!setptrs!},
\passthrough{\lstinline!levwk*!} \\
50/4 & Bound-bound radiative line & wavelength/energy, oscillator
strength, A value, upper/lower level indices &
\passthrough{\lstinline!ucalc!} type 50,
\passthrough{\lstinline!calc\_hmc\_ion!} \\
53/7 & OP bound-free photoionization & threshold, energy-offset grid,
cross-section grid, destination levels & \passthrough{\lstinline!ucalc!}
type 53, \passthrough{\lstinline!phint53!} \\
56/3 & Tabulated effective collision strengths & temperature grid,
upsilon values, level pair & \passthrough{\lstinline!upsil*!}, collision
decoder \\
63/3 & Bautista hydrogenic collisions and same-n l mixing & quantum
numbers, fit/control integers, level pair &
\passthrough{\lstinline!amcrs!}, \passthrough{\lstinline!anl1!},
\passthrough{\lstinline!erc!}, \passthrough{\lstinline!velimp!} \\
68/3 & He-like collision strengths & Zhang-Sampson parameters and
transition indices & \passthrough{\lstinline!calt68!} \\
69/3 & He-like fitted collision strengths & Kato-Nakazaki fit parameters
and transition indices & \passthrough{\lstinline!calt69!},
\passthrough{\lstinline!calc\_kato!} \\
70/7 & Superlevel photoionization & superlevel cross-section
coefficients and parent mapping & \passthrough{\lstinline!calt70!},
\passthrough{\lstinline!phint53hunt!} \\
71/14 & Superlevel radiative cascade & A value and
superlevel-to-spectroscopic destination &
\passthrough{\lstinline!calt71!} \\
74/7 & Dielectronic-recombination resonance contribution &
delta-function resonance strength and destination mapping &
\passthrough{\lstinline!calt74!} \\
77/23 & Superlevel collisional cascade & density-scaled collisional
partner of type 71 & \passthrough{\lstinline!calt77!} \\
99/7 & Superlevel source/sink closure & recombination normalization plus
photoionization search & \passthrough{\lstinline!calt99!},
\passthrough{\lstinline!phint53hunt!} \\
\end{longtable}

\subsection{4. XSTAR source-code layout}\label{xstar-source-code-layout}

\subsubsection{4.1 Top-level directories}\label{top-level-directories}

\begin{longtable}[]{@{}
  >{\raggedright\arraybackslash}p{(\columnwidth - 2\tabcolsep) * \real{0.5000}}
  >{\raggedright\arraybackslash}p{(\columnwidth - 2\tabcolsep) * \real{0.5000}}@{}}
\toprule\noalign{}
\begin{minipage}[b]{\linewidth}\raggedright
Path
\end{minipage} & \begin{minipage}[b]{\linewidth}\raggedright
Role
\end{minipage} \\
\midrule\noalign{}
\endhead
\bottomrule\noalign{}
\endlastfoot
\passthrough{\lstinline!src/xstar/xstar.f90!} & Executable driver, input
handling, radial shell/pass loop, final products. \\
\passthrough{\lstinline!xstarlib/src/!} & Atomic data, local plasma
solve, transfer, emissivity, matrix, numerical routines. \\
\passthrough{\lstinline!data/!} & Runtime auxiliary data; the large
\passthrough{\lstinline!atdb.fits!} is normally distributed/installed
separately. \\
\passthrough{\lstinline!manual/!} and \passthrough{\lstinline!doc/!} &
XSTAR manual and supporting documentation. \\
\passthrough{\lstinline!utils/!}, \passthrough{\lstinline!scripts/!} &
Conversion, setup, and operational helpers. \\
\end{longtable}

\subsubsection{4.2 Functional source map}\label{functional-source-map}

\begin{longtable}[]{@{}
  >{\raggedright\arraybackslash}p{(\columnwidth - 4\tabcolsep) * \real{0.3333}}
  >{\raggedright\arraybackslash}p{(\columnwidth - 4\tabcolsep) * \real{0.3333}}
  >{\raggedright\arraybackslash}p{(\columnwidth - 4\tabcolsep) * \real{0.3333}}@{}}
\toprule\noalign{}
\begin{minipage}[b]{\linewidth}\raggedright
Stage
\end{minipage} & \begin{minipage}[b]{\linewidth}\raggedright
Principal routines
\end{minipage} & \begin{minipage}[b]{\linewidth}\raggedright
Live data produced/consumed
\end{minipage} \\
\midrule\noalign{}
\endhead
\bottomrule\noalign{}
\endlastfoot
Database read & \passthrough{\lstinline!readtbl!},
\passthrough{\lstinline!dread!}, \passthrough{\lstinline!setptrs!},
\passthrough{\lstinline!globaldata!} & \passthrough{\lstinline!nptrs!},
\passthrough{\lstinline!rdat1!}, \passthrough{\lstinline!idat1!},
\passthrough{\lstinline!kdat1!}, derived pointers. \\
Continuum setup & \passthrough{\lstinline!ener!}, incident-spectrum
routines & high-resolution \passthrough{\lstinline!epi!},
incident/diffuse spectral arrays. \\
Transfer & \passthrough{\lstinline!trnfrc!},
\passthrough{\lstinline!trnfrn!}, \passthrough{\lstinline!bremsmap!} &
\passthrough{\lstinline!bremsa!}, \passthrough{\lstinline!bremsint!},
mapped \passthrough{\lstinline!epim!},
\passthrough{\lstinline!bremsam!}. \\
Thermal/charge iteration & \passthrough{\lstinline!dsec!},
\passthrough{\lstinline!heatf!}, \passthrough{\lstinline!heatt!} &
temperature, electron fraction, heating/cooling. \\
Ion rates & \passthrough{\lstinline!calc\_ion\_rates!},
\passthrough{\lstinline!istruc!}, \passthrough{\lstinline!ioneqm!} &
total ionization/recombination rates and ion fractions. \\
Element basis & \passthrough{\lstinline!calc\_hmc\_all!},
\passthrough{\lstinline!calc\_hmc\_element!},
\passthrough{\lstinline!levwkelement!} & compact
\passthrough{\lstinline!ipmat2!} basis, ion/superlevel maps, initial
populations. \\
Record rates & \passthrough{\lstinline!calc\_hmc\_ion!},
\passthrough{\lstinline!ucalc!}, \passthrough{\lstinline!calt*!},
\passthrough{\lstinline!phint53*!} &
\passthrough{\lstinline!ans1..ans6!}, rate provenance,
opacity/emissivity contributions. \\
Matrix solve & \passthrough{\lstinline!msolvelucy!},
\passthrough{\lstinline!leqt2f!}, LU helpers & level/superlevel
population vector \passthrough{\lstinline!x!}. \\
Emissivity/opacity & \passthrough{\lstinline!calc\_emis*!},
\passthrough{\lstinline!calc\_emisab*!},
\passthrough{\lstinline!linopac!} & lines, RRCs, continuum emissivity
and opacity. \\
Depth update & \passthrough{\lstinline!stpcut!} & direction-dependent
\passthrough{\lstinline!tau0!}, \passthrough{\lstinline!tauc!},
continuum depths. \\
Output & \passthrough{\lstinline!savd!},
\passthrough{\lstinline!fstepr*!}, \passthrough{\lstinline!pprint!},
\passthrough{\lstinline!writespectra*!} &
\passthrough{\lstinline!xoNN\_detail/detal*!} and
\passthrough{\lstinline!xout\_*!} products. \\
\end{longtable}

\subsubsection{4.3 Output products and their debugging
value}\label{output-products-and-their-debugging-value}

\begin{longtable}[]{@{}
  >{\raggedright\arraybackslash}p{(\columnwidth - 4\tabcolsep) * \real{0.3333}}
  >{\raggedright\arraybackslash}p{(\columnwidth - 4\tabcolsep) * \real{0.3333}}
  >{\raggedright\arraybackslash}p{(\columnwidth - 4\tabcolsep) * \real{0.3333}}@{}}
\toprule\noalign{}
\begin{minipage}[b]{\linewidth}\raggedright
Product
\end{minipage} & \begin{minipage}[b]{\linewidth}\raggedright
Main contents
\end{minipage} & \begin{minipage}[b]{\linewidth}\raggedright
Parity use
\end{minipage} \\
\midrule\noalign{}
\endhead
\bottomrule\noalign{}
\endlastfoot
\passthrough{\lstinline!xout\_abund1.fits!} & zone temperature, electron
fraction/density information, ion fractions & local-state selection and
ionization closure. \\
\passthrough{\lstinline!xout\_lines1.fits!} & final line
luminosities/emissivities & f/i/r benchmark target, not a live rate
field. \\
\passthrough{\lstinline!xout\_rrc1.fits!},
\passthrough{\lstinline!xout\_cont1.fits!},
\passthrough{\lstinline!xout\_spect1.fits!} & RRC/continuum/final
spectrum & output-layer validation. \\
\passthrough{\lstinline!xoNN\_detail.fits!} & level populations by zone
& post-solve population parity. \\
\passthrough{\lstinline!xoNN\_detal2.fits!} & line emissivity, opacity,
inward/outward optical depths & type-50 escape and line-transfer
audits. \\
\passthrough{\lstinline!xoNN\_detal3.fits!} & RRC emissivity, opacity,
depths & bound-free escape audits. \\
\passthrough{\lstinline!xoNN\_detal4.fits!} & continuum
emissivity/opacity/depth arrays & reconstruction aid, but not identical
to live \passthrough{\lstinline!bremsam!}. \\
\end{longtable}

\subsection{5. From an ATDB record to a matrix
coefficient}\label{from-an-atdb-record-to-a-matrix-coefficient}

The record-to-solution path is:

\begin{lstlisting}
ATDB record number
  -> POINTERS header
  -> REALS/INTEGERS/CHARS payload slices
  -> setptrs parent and rate-type chains
  -> calc_hmc_ion chooses local destinations and escape context
  -> ucalc dispatches by data type
  -> ans1..ans6 raw physical outputs
  -> possible branch-specific swaps/sign changes
  -> four ajisi/indbi insertions for a two-way transition
  -> calc_hmc_element adds ion-block ipmat2 offsets
  -> msolvelucy solves the compact element system
\end{lstlisting}

For a matrix-active two-way process,
\passthrough{\lstinline!calc\_hmc\_ion!} typically emits four entries:
two off-diagonal gains and two diagonal losses. A rate comparison is
incomplete unless both the numeric values and these row/column
placements agree.

\subsubsection{5.1 Compact element basis and parent-continuum
aliases}\label{compact-element-basis-and-parent-continuum-aliases}

XSTAR does not concatenate independent ion level lists naively. Within
\passthrough{\lstinline!calc\_hmc\_element!}, the compact basis advances
approximately as:

\begin{lstlisting}[language=Fortran]
ipmat2 = ipmat2 + nlev - 1
\end{lstlisting}

and adds a final parent/continuum slot after the ion loop. Consequently,
the last row of one ion block can be the same compact unknown as the
ground/parent role of the adjacent ion. One compact
\passthrough{\lstinline!ipmat2!} row may therefore have multiple
physical role labels.

The v0.3.192--v0.3.198 probes established the exact O-element topology
for the selected O VII benchmark:

\begin{itemize}
\tightlist
\item
  607 compact XSTAR population rows;
\item
  six ion blocks in the selected element solve;
\item
  six shared alias rows;
\item
  all matrix endpoints mapped by
  \passthrough{\lstinline!compact\_ipmat2 = ion\_ipmat2\_offset + indbi!};
\item
  all unrelated-element records filtered before compact mapping.
\end{itemize}

\subsection{6. Physical equations and their source-code
roles}\label{physical-equations-and-their-source-code-roles}

\subsubsection{6.1 Radiation and ionization
parameter}\label{radiation-and-ionization-parameter}

XSTAR commonly characterizes the incident radiation by

\[
\xi = \frac{L_{\mathrm{ion}}}{n\,r^2}
\quad [\mathrm{erg\,cm\,s^{-1}}],
\]

where the exact density convention follows the run definition. The local
rate calculation does not use only \(\xi\); it uses the transferred
spectral field. In the current source path:

\begin{itemize}
\tightlist
\item
  \passthrough{\lstinline!epi(:)!} is an energy grid;
\item
  \passthrough{\lstinline!bremsa(:)!} is the live high-resolution
  ionizing field used by transfer;
\item
  \passthrough{\lstinline!epim(:)!},
  \passthrough{\lstinline!bremsam(:)!}, and
  \passthrough{\lstinline!bremsint(:)!} are the mapped rate grid passed
  to atomic calculations.
\end{itemize}

\passthrough{\lstinline!bremsmap.f90!} primarily maps the
high-resolution field to the rate grid by bin lookup and constructs its
cumulative integral. It does not contain a hidden factor near 44.

\subsubsection{6.2 Atomic levels and
wavelengths}\label{atomic-levels-and-wavelengths}

For levels \(i\) and \(j\),

\[
\Delta E_{ij} = E_j-E_i, \qquad
\lambda_{ji} = \frac{hc}{|\Delta E_{ij}|}.
\]

Each level also carries a statistical weight \(g_i\). XSTAR uses level
energies and weights for transition ordering, detailed balance, LTE
ratios, continuum thresholds, and matrix destinations. The database
wavelength may be retained for line identification, while energy
differences can be used for energy conservation in emissivity
calculations.

\subsubsection{6.3 Statistical equilibrium of level
populations}\label{statistical-equilibrium-of-level-populations}

For each compact population row \(i\), steady state requires

\[
0 = \sum_{j\ne i} n_j R_{j\rightarrow i}
    -n_i\sum_{j\ne i}R_{i\rightarrow j}
    +S_i,
\]

where \(S_i\) includes processes represented through adjacent-ion and
source closure. In matrix form,

\[
\mathbf{{A}}\,\mathbf{{x}}=\mathbf{{b}},
\]

with one row replaced by a normalization condition such as

\[
\sum_i x_i=1.
\]

\passthrough{\lstinline!calc\_hmc\_ion!} creates sparse transition
entries, \passthrough{\lstinline!calc\_hmc\_element!} maps them into
compact element coordinates, and \passthrough{\lstinline!msolvelucy!}
solves through a superlevel condensation/fixed-point procedure. The
row-balance audit evaluates

\[
r_i=\sum_j A_{{ij}}x_j
\]

against the captured post-\passthrough{\lstinline!msolvelucy!} vector.

\subsubsection{6.4 Ionization equilibrium}\label{ionization-equilibrium}

A simplified adjacent-stage balance is

\[
n_q\,\Gamma_q = n_{{q+1}}\,n_e\,\alpha_{{q+1}},
\]

but XSTAR assembles all included stages and channels. In a general
steady-state form,

\[
0=\sum_{{q'\ne q}} n_{{q'}} I_{{q'\rightarrow q}}
  -n_q\sum_{{q'\ne q}} I_{{q\rightarrow q'}},
\qquad
\sum_q f_q=1.
\]

\passthrough{\lstinline!calc\_ion\_rates!} accumulates total
photoionization and recombination terms;
\passthrough{\lstinline!istruc!}/\passthrough{\lstinline!ioneqm!} solve
the ion fractions. Level-resolved and superlevel source terms then
couple the ion solution to the element population matrix.

\subsubsection{6.5 Photoionization and heating: type
53}\label{photoionization-and-heating-type-53}

For a level with threshold \(E_0\), the physical forward rate is
represented by

\[
\Gamma_{{\rm PI}} = \int_{{E_0}}^\infty
\sigma(E)\,\frac{F_E}{E}\,dE,
\]

and the corresponding photoelectron heating is

\[
H_{{\rm PI}} = n_i\int_{{E_0}}^\infty
(E-E_0)\,\sigma(E)\,\frac{F_E}{E}\,dE.
\]

\passthrough{\lstinline!phint53.f90!} maps the OP cross-section grid
onto the XSTAR energy grid and trapezoid-integrates the rate and heating
terms. Its code contains paired factors of 12.56 that cancel in the
forward photoionization integrand. The reverse branch uses a
Milne/detailed-balance construction, statistical-weight/LTE factors, and
the local temperature.

\paragraph{The unresolved approximately 44
scale}\label{the-unresolved-approximately-44-scale}

The instrumented live-rate-grid audit found

\begin{lstlisting}
preserved Python type-53 proxy matrix rate / XSTAR live phint53 ans1 ≈ 44.2
\end{lstlisting}

This is not a hidden \passthrough{\lstinline!bremsmap!} or
\passthrough{\lstinline!phint53!} constant. It is currently interpreted
as the ratio of the legacy diagnostic
\passthrough{\lstinline!xstar-powerlaw!} normalization to the
source-code-equivalent live \passthrough{\lstinline!bremsam!} rate.
Replaying exact XSTAR \passthrough{\lstinline!ucalc!} type-53 rates
changed the O VII triplet population fractions by only about
\(8.4\times10^{{-7}}\), so this scale is real but was not the dominant
cause of the earlier triplet population mismatch.

The planned fix is not an empirical division by 44. The production path
must:

\begin{enumerate}
\def\labelenumi{\arabic{enumi}.}
\tightlist
\item
  carry the exact local \passthrough{\lstinline!epim!},
  \passthrough{\lstinline!bremsam!}, and
  \passthrough{\lstinline!bremsint!} state;
\item
  port the complete type-53 threshold mapping and
  \passthrough{\lstinline!phint53!} integral;
\item
  compare native Python \passthrough{\lstinline!ans1..ans6!} with the
  same XSTAR \passthrough{\lstinline!ucalc!} capture;
\item
  replace proxy matrix terms only after branch-level and
  matrix-placement parity pass.
\end{enumerate}

\subsubsection{6.6 Recombination and the Milne
relation}\label{recombination-and-the-milne-relation}

Radiative recombination is the inverse of photoionization.
Schematically,

\[
\alpha_i(T) = \int_0^\infty v\,f_M(v,T)\,
\sigma_{{\rm RR},i}(v)\,dv,
\]

where the recombination cross section is related to the photoionization
cross section by detailed balance. XSTAR's implementation uses
statistical-weight factors, an LTE population ratio, the Maxwell factor
\(\exp[-(E-E_0)/kT]\), and the same cross-section grid. The
\passthrough{\lstinline!phint53!} reverse side also generates RRC
emissivity and recombination cooling.

Total recombination data types provide fitted rates, while types 70, 74,
and 99 participate in explicit superlevel and dielectronic-recombination
closure. A full implementation must not replace these with arbitrary
source vectors.

\subsubsection{6.7 Electron-impact excitation and
de-excitation}\label{electron-impact-excitation-and-de-excitation}

For an effective collision strength \(\Upsilon_{{ij}}(T)\), the common
Maxwellian form is

\[
q_{{ij}} = \frac{{8.629\times10^{{-6}}}}{{g_i T^{{1/2}}}}
\Upsilon_{{ij}}(T)\exp\left(-\frac{{\Delta E_{{ij}}}}{{kT}}\right)
\quad \mathrm{{cm^3\,s^{{-1}}}},
\]

\[
q_{{ji}} = \frac{{8.629\times10^{{-6}}}}{{g_j T^{{1/2}}}}
\Upsilon_{{ij}}(T).
\]

Matrix rates are \(C_{{ij}}=n_e q_{{ij}}\). In the current He-like work:

\begin{itemize}
\tightlist
\item
  type 56 carries tabulated effective collision strengths;
\item
  type 63 implements Bautista hydrogenic transitions and same-\(n\)
  \(l\)-mixing branches;
\item
  type 68 uses Zhang-Sampson He-like collision strengths;
\item
  type 69 uses Kato-Nakazaki fits;
\item
  type 77 provides superlevel-to-spectroscopic collisional coupling.
\end{itemize}

Direct probes established near-unity record/matrix parity for types 63,
68, and 69 in the selected O VII local state.

\subsubsection{6.8 Radiative decay, escape, and line pumping: type
50}\label{radiative-decay-escape-and-line-pumping-type-50}

For a line with spontaneous rate \(A_{{ul}}\), XSTAR first constructs
direction/covering factors

\[
p_1=P_{{\rm esc}}(\tau_{{\rm in}})(1-C_f),
\]

\[
p_2=P_{{\rm esc}}(\tau_{{\rm out}})(1-C_f)
 +2P_{{\rm esc}}(\tau_{{\rm in}}+\tau_{{\rm out}})C_f.
\]

The escaped decay rate before the final branch swap is

\[
R_{{u\rightarrow l}}^{{\rm esc}}=A_{{ul}}(p_1+p_2).
\]

The line-center velocity-integrated cross-section factor is represented
in the source as

\[
\sigma_v = 0.02655\,f_{{lu}}\,\frac{{\lambda_{{\rm cm}}}}{{v_{{\rm th/turb}}}},
\]

and the photoexcitation/pumping rate is approximately

\[
R_{{l\rightarrow u}} = \sigma_v\,
\mathrm{{bremsa}}(E_{{ul}})\,\frac{{v_{{\rm th/turb}}}}{{c}}
\,f_{{\rm linabs}}\,(1-C_f).
\]

\passthrough{\lstinline!ucalc!} finally swaps
\passthrough{\lstinline!ans1!} and \passthrough{\lstinline!ans2!} so the
returned forward branch is excitation and the reverse branch is escaped
decay. For \(C_f=1\), pumping is suppressed and the escape term uses the
combined inward+outward optical depth. The v0.3.156 detail-state audit
verified all five O VII triplet type-50 rows and the matrix insertions.

\subsubsection{6.9 Line width, opacity, and
emissivity}\label{line-width-opacity-and-emissivity}

\passthrough{\lstinline!linopac.f90!} combines thermal and turbulent
widths in quadrature:

\[
\Delta E_D = \sqrt{{\Delta E_{{\rm th}}^2+\Delta E_{{\rm turb}}^2}},
\]

with \(v_{{\rm th}}\propto\sqrt{{T/A}}\) and \(\Delta E/E=v/c\). It
deposits a Gaussian/Voigt line profile into the continuum opacity grid
or uses a single-bin approximation in fast mode.

For a radiative transition, the optically modified emissivity is
proportional to

\[
\epsilon_{{ul}} = n_u A_{{ul}} h\nu\,P_{{\rm esc}},
\]

with inward/outward pieces separated by the escape factors. Net line
output may also subtract radiative excitation/absorption, as in the
\passthrough{\lstinline!calc\_emis\_ion!} expression involving
\passthrough{\lstinline!ans2*abund2 - ans1*abund1!}.

\subsubsection{6.10 Superlevels and compact
closure}\label{superlevels-and-compact-closure}

Superlevels reduce a large level system by grouping rows through
\passthrough{\lstinline!nsup!}. \passthrough{\lstinline!msolvelucy!}
computes within-superlevel fractions

\[
r_i=\frac{{x_i}}{{p_{{s(i)}}}},
\qquad
p_s=\sum_{{i\in s}}x_i,
\]

constructs a condensed superlevel matrix, imposes number conservation,
solves for \(p_s\), and reconstructs \(x_i=r_i p_{{s(i)}}\). Fixed-point
iterations then update level fractions and superlevel totals.

Types 71 and 77 connect superlevels to spectroscopic levels radiatively
and collisionally. Types 70, 74, and 99 couple photoionization,
dielectronic recombination, parent continua, and superlevel source/sink
closure. These branches are essential for a physical full-element
solution.

\subsubsection{6.11 Thermal balance}\label{thermal-balance}

The local temperature is determined by

\[
H(T,n_e,\{{f_q\}},J_E)-C(T,n_e,\{{f_q\}},J_E)=0,
\]

with charge/electron consistency solved simultaneously or iteratively.
Heating includes photoelectrons, Compton and other channels; cooling
includes line emission, recombination, free-free emission, and other
atomic channels. \passthrough{\lstinline!dsec!} repeatedly calls the
population/ionization calculation while adjusting temperature and
electron fraction.

\subsubsection{6.12 Radial transfer and optical-depth
evolution}\label{radial-transfer-and-optical-depth-evolution}

The transferred field at a zone depends on incident radiation,
attenuation, and diffuse emission. XSTAR uses direction-dependent
continuum/line depths and a single-stream approximation in the current
transfer path. \passthrough{\lstinline!trnfrc!} constructs the local
continuum used by the rates; \passthrough{\lstinline!stpcut!} updates
line and continuum optical depths after the local solution. A full
replacement must therefore iterate transfer and plasma state rather than
solve isolated zones only.

\subsection{7. He-like C V, O VII, Mg XI, and Ca XIX
calculations}\label{he-like-c-v-o-vii-mg-xi-and-ca-xix-calculations}

\subsubsection{7.1 Triplet diagnostics}\label{triplet-diagnostics}

The principal He-like lines are:

\begin{itemize}
\tightlist
\item
  forbidden \(f\): \(1s2s\,^3S_1\rightarrow1s^2\,^1S_0\);
\item
  intercombination \(i\): \(1s2p\,^3P_{{1,2}}\rightarrow1s^2\,^1S_0\),
  often summed;
\item
  resonance \(r\): \(1s2p\,^1P_1\rightarrow1s^2\,^1S_0\).
\end{itemize}

The standard diagnostics are

\[
R=\frac{{f}}{{i}},
\qquad
G=\frac{{f+i}}{{r}}.
\]

\(R\) responds strongly to density and radiation-driven transfer between
\(^3S\) and \(^3P\); \(G\) responds to the balance of
recombination/cascade feeding, collision excitation, and resonance-line
transfer.

\subsubsection{7.2 Canonical same-run benchmark
suite}\label{canonical-same-run-benchmark-suite}

\begin{longtable}[]{@{}lrrl@{}}
\toprule\noalign{}
Ion & Atomic number & Default wavelength window & Run directory \\
\midrule\noalign{}
\endhead
\bottomrule\noalign{}
\endlastfoot
C V & 6 & 40.0--42.0 Å &
\passthrough{\lstinline!helike\_type69/c5\_ne1e8!} \\
O VII & 8 & 21.0--23.0 Å &
\passthrough{\lstinline!helike\_type69/o7\_ne1e8!} \\
Mg XI & 12 & 9.0--9.4 Å &
\passthrough{\lstinline!helike\_type69/mg11\_ne1e8!} \\
Ca XIX & 20 & 3.0--3.35 Å &
\passthrough{\lstinline!helike\_type69/ca19\_xi3\_ne1e8!} \\
\end{longtable}

These cases test the same atomic architecture across increasing nuclear
charge and line energy. The exact local temperature, electron density,
ion fraction, and \(\log\xi\) are extracted from each same-run XSTAR
output rather than imposed from one universal model.

\subsubsection{7.3 Physical processes that control the four-ion
benchmark}\label{physical-processes-that-control-the-four-ion-benchmark}

\begin{longtable}[]{@{}
  >{\raggedright\arraybackslash}p{(\columnwidth - 4\tabcolsep) * \real{0.3333}}
  >{\raggedright\arraybackslash}p{(\columnwidth - 4\tabcolsep) * \real{0.3333}}
  >{\raggedright\arraybackslash}p{(\columnwidth - 4\tabcolsep) * \real{0.3333}}@{}}
\toprule\noalign{}
\begin{minipage}[b]{\linewidth}\raggedright
Process
\end{minipage} & \begin{minipage}[b]{\linewidth}\raggedright
Main data/source branches
\end{minipage} & \begin{minipage}[b]{\linewidth}\raggedright
Effect on f/i/r
\end{minipage} \\
\midrule\noalign{}
\endhead
\bottomrule\noalign{}
\endlastfoot
Spontaneous decay and escape & type 50, \passthrough{\lstinline!pescl!},
\passthrough{\lstinline!calc\_hmc\_ion!} & Direct line drain; optical
depth changes apparent resonance strength. \\
Line pumping & type 50 plus live radiation and
\passthrough{\lstinline!cfrac!} & Transfers population upward;
suppressed for \passthrough{\lstinline!cfrac=1!}. \\
Direct electron excitation & types 63, 68, 69 & Feeds triplet and
resonance upper levels with strong temperature/density dependence. \\
Radiative cascades & type 71 & Feeds spectroscopic levels from
superlevels. \\
Collisional superlevel coupling & type 77 & Redistributes superlevel
populations. \\
Photoionization/RR & type 53 & Couples levels to parent continuum and
produces RRC/heating/cooling. \\
DR/source closure & types 70, 74, 99 & Supplies and drains
superlevels/parent continua. \\
Adjacent-ion compact aliases &
\passthrough{\lstinline!calc\_hmc\_element!} topology & Ensures
ion-stage and level-population normalization are solved consistently. \\
\end{longtable}

\subsubsection{7.4 Current benchmark
interpretation}\label{current-benchmark-interpretation}

O VII is the deepest source-code parity case because it has direct
rate-grid, \passthrough{\lstinline!ucalc!}, matrix, population, basis,
and balance probes. C V, Mg XI, and Ca XIX are retained as cross-ion
regression cases. The aim is not to tune each ion separately; it is to
implement one source-equivalent data/rate/basis pipeline that reproduces
all four under their own local XSTAR states.

\subsection{8. New findings established by the validation
sequence}\label{new-findings-established-by-the-validation-sequence}

\subsubsection{8.1 Type-50 escape and
pumping}\label{type-50-escape-and-pumping}

\begin{itemize}
\tightlist
\item
  The exact \passthrough{\lstinline!cfrac!}-dependent
  \passthrough{\lstinline!ptmp1/ptmp2!} formulas are required.
\item
  For the O VII benchmark with \passthrough{\lstinline!cfrac=1!}, line
  pumping is zero.
\item
  Five selected O VII triplet type-50 records match
  \passthrough{\lstinline!ucalc!} and matrix insertions.
\item
  A global diagnostic escape factor such as 0.35 must not be used as the
  production source-equivalent treatment.
\end{itemize}

\subsubsection{8.2 Matrix-family parity}\label{matrix-family-parity}

At the selected O VII local state:

\begin{itemize}
\tightlist
\item
  type 63, 68, and 69 collision matrix terms are essentially at unity
  parity;
\item
  type 71 cascade terms are at approximately 0.998 median parity;
\item
  type 50 contains exact triplet matches but older non-triplet proxy
  treatments remain in historical products;
\item
  type 77 showed a moderate historical difference and remains a
  native-port target;
\item
  type 53 historical proxy rows differ strongly from live
  \passthrough{\lstinline!phint53!} rates.
\end{itemize}

\subsubsection{8.3 The type-53 scale is real but not the main population
lever}\label{the-type-53-scale-is-real-but-not-the-main-population-lever}

Live \passthrough{\lstinline!epim/bremsam!} recomputation confirmed the
approximately 44 scale. Controlled replay of exact XSTAR
\passthrough{\lstinline!ucalc!} rates for type 53, 50, 77, and 99 moved
O VII triplet population fractions by at most about
\(8.36\times10^{{-7}}\). This established that population/source closure
and basis mapping had to be fixed before interpreting rate-family
changes end to end.

\subsubsection{8.4 Direct Fortran probes}\label{direct-fortran-probes}

The instrumentation now captures:

\begin{itemize}
\tightlist
\item
  per-record \passthrough{\lstinline!ucalc ans1..ans6!};
\item
  the four \passthrough{\lstinline!calc\_hmc\_ion!} matrix insertions
  and shared capture IDs;
\item
  live \passthrough{\lstinline!epim!},
  \passthrough{\lstinline!bremsam!}, and
  \passthrough{\lstinline!bremsint!} rate grids;
\item
  paired before/after-\passthrough{\lstinline!msolvelucy!} population
  vectors;
\item
  exact \passthrough{\lstinline!calc\_hmc\_element!} compact-basis roles
  and ion-block offsets.
\end{itemize}

These probes convert implementation questions into direct record/row
comparisons rather than inferred fits to final line ratios.

\subsubsection{8.5 Exact O-element compact
basis}\label{exact-o-element-compact-basis}

The selected O-element solve has:

\begin{itemize}
\tightlist
\item
  607 compact population rows;
\item
  352 nonzero post-solve rows;
\item
  six compact ion blocks and shared parent-continuum aliases.
\end{itemize}

The original sequential Python
\passthrough{\lstinline!xstar\_ipmat2\_index!} mapping was wrong.
Physical remapping by
\passthrough{\lstinline!(ion\_stage, level\_index)!} maps all 114
current Python population identities to 113 unique compact rows and
raises captured XSTAR population coverage from about
\(2.6\times10^{{-6}}\) to

\[
0.991904690445624.
\]

\subsubsection{8.6 Six-row priority
expansion}\label{six-row-priority-expansion}

Adding compact rows

\begin{lstlisting}
293, 241, 244, 242, 80, 79
\end{lstlisting}

raises represented solved-population coverage to

\[
0.9999999463123856.
\]

Three are ordinary rows and three are shared
parent-continuum/next-ion-ground aliases. The active staged basis
contains 119 unique compact rows represented by 120 physical Python
identities.

\subsubsection{8.7 Complete priority-row matrix manifest and
balance}\label{complete-priority-row-matrix-manifest-and-balance}

After correcting endpoint translation and element filtering:

\begin{itemize}
\tightlist
\item
  all six priority rows have Fortran matrix coverage;
\item
  1,260 matrix-active records have four rows;
\item
  nine intentional non-matrix metadata records have zero rows;
\item
  no compact matrix endpoints are unmapped;
\item
  the six selected row equations pass a 0.5\% relative residual
  threshold against the captured XSTAR population vector;
\item
  the maximum row residual is approximately \(1.7471\times10^{{-3}}\).
\end{itemize}

This is the strongest current evidence that the reconstructed topology
and probed terms are correct for the staged subset.

\subsection{\texorpdfstring{9. \texttt{xstar-atomic} data and software
structures}{9. xstar-atomic data and software structures}}\label{xstar-atomic-data-and-software-structures}

\begin{longtable}[]{@{}
  >{\raggedright\arraybackslash}p{(\columnwidth - 2\tabcolsep) * \real{0.5000}}
  >{\raggedright\arraybackslash}p{(\columnwidth - 2\tabcolsep) * \real{0.5000}}@{}}
\toprule\noalign{}
\begin{minipage}[b]{\linewidth}\raggedright
Python structure/module
\end{minipage} & \begin{minipage}[b]{\linewidth}\raggedright
Purpose
\end{minipage} \\
\midrule\noalign{}
\endhead
\bottomrule\noalign{}
\endlastfoot
\passthrough{\lstinline!ATDB!}, \passthrough{\lstinline!RecordHeader!},
\passthrough{\lstinline!RawRecord!} & Low-level packed FITS access. \\
\passthrough{\lstinline!ATDBIndexArrays!},
\passthrough{\lstinline!IndexedRecord!},
\passthrough{\lstinline!ElementInfo!}, \passthrough{\lstinline!IonInfo!}
& Array-backed element/ion/level hierarchy and targeted selection. \\
\passthrough{\lstinline!XSTARAtomic!} & Reusable high-level database
object and workflow namespaces. \\
\passthrough{\lstinline!LocalPlasmaState!} & Temperature, density,
ionization and local abundance state. \\
\passthrough{\lstinline!RadiationField!} & Energy grid and local
radiation samples; future source-equivalent rate-grid carrier. \\
\passthrough{\lstinline!EscapeContext!} &
\passthrough{\lstinline!cfrac!}, line/continuum optical depths,
velocity/escape data. \\
\passthrough{\lstinline!XSTARContext!} & Combined state/radiation/escape
object. \\
\passthrough{\lstinline!RateEvaluation!} & Numeric rate plus
raw/interpreted branch provenance. \\
\passthrough{\lstinline!XSTARContinuumState!},
\passthrough{\lstinline!XSTARZoneState!},
\passthrough{\lstinline!XSTARRunState!} & Skeleton for future full
radial calculations and output recreation. \\
Record-level parity products & Join Python records to XSTAR
\passthrough{\lstinline!ucalc!} and four matrix rows. \\
Element-basis probe/remap/scaffold & Reconstructs compact
\passthrough{\lstinline!ipmat2!} topology and aliases. \\
Priority expansion/closure/balance products & Staged native
implementation manifest and consistency gates. \\
\end{longtable}

\subsection{10. Current implementation
status}\label{current-implementation-status}

\subsubsection{Source-equivalent or strongly
validated}\label{source-equivalent-or-strongly-validated}

\begin{itemize}
\tightlist
\item
  packed FITS vector reader and hierarchy;
\item
  level and line extraction;
\item
  selected photoionization and recombination decoders;
\item
  type 56 and validated type 63 collision branches;
\item
  type 68 and 69 He-like collision evaluation/matrix placement;
\item
  type 50 triplet detail-state escape/pumping calculation;
\item
  type 71 cascade evaluation/matrix placement;
\item
  raw record-level \passthrough{\lstinline!ucalc!} and matrix probe
  association;
\item
  exact O-element compact-basis reconstruction and physical remapping;
\item
  six-row priority matrix manifest and row-balance validation.
\end{itemize}

\subsubsection{Diagnostic or incomplete}\label{diagnostic-or-incomplete}

\begin{itemize}
\tightlist
\item
  historical proxy \passthrough{\lstinline!xstar-powerlaw!} type-53
  matrix normalization;
\item
  native type 77 implementation over the expanded basis;
\item
  native type 70/74/99 parent/superlevel source closure;
\item
  expanded-basis RHS and normalization assembly;
\item
  complete 607-row native element solve;
\item
  coupled ionization/temperature iteration;
\item
  radial transfer and output reproduction.
\end{itemize}

\subsection{11. Roadmap to a full XSTAR
implementation}\label{roadmap-to-a-full-xstar-implementation}

\subsubsection{Phase A: controlled 119-row compact
solve}\label{phase-a-controlled-119-row-compact-solve}

\begin{enumerate}
\def\labelenumi{\arabic{enumi}.}
\tightlist
\item
  Instantiate the corrected 113 mapped compact rows plus the six
  priority rows.
\item
  Preserve all parent-continuum aliases as one unknown with multiple
  role labels.
\item
  Assemble probed matrix entries and explicit RHS/normalization terms.
\item
  Solve the 119-row system and compare every population and selected row
  residual with XSTAR.
\item
  Keep this as a replay/parity mode until native formulas reproduce it.
\end{enumerate}

\textbf{Acceptance gates:} no unmapped endpoints; normalization exactly
one; priority rows below the residual tolerance; mapped population
fractions agree within declared tolerances.

\subsubsection{Phase B: native rates for the priority
subset}\label{phase-b-native-rates-for-the-priority-subset}

Port only the rate families actually touching the six selected rows.
Every native branch must reproduce:

\begin{lstlisting}
raw ATDB payload -> ans1..ans6 -> four matrix rows -> row balance
\end{lstlisting}

No empirical source weights or scale factors should be accepted in
production mode.

\subsubsection{Phase C: eliminate the type-53 approximately 44
scale}\label{phase-c-eliminate-the-type-53-approximately-44-scale}

\begin{enumerate}
\def\labelenumi{\arabic{enumi}.}
\tightlist
\item
  Promote \passthrough{\lstinline!epim!},
  \passthrough{\lstinline!bremsam!}, \passthrough{\lstinline!bremsint!}
  to first-class live-state arrays.
\item
  Port complete \passthrough{\lstinline!phint53!} and
  \passthrough{\lstinline!phint53hunt!} behavior.
\item
  Validate threshold binning, cross-section mapping, forward rate,
  reverse rate, heating/cooling, opacity, and RRC emissivity.
\item
  Replace proxy type-53 terms only after record-level parity.
\item
  Repeat across C V, O VII, Mg XI, and Ca XIX.
\end{enumerate}

\subsubsection{Phase D: expand from 119 to all 607 compact
rows}\label{phase-d-expand-from-119-to-all-607-compact-rows}

Use the population-ranked scaffold to add remaining rows in tiers. For
each tier:

\begin{itemize}
\tightlist
\item
  add row identities and aliases;
\item
  add all touching records and endpoints;
\item
  verify row balance;
\item
  compare population coverage and line emissivities;
\item
  retain a deterministic mapping from physical roles to compact rows.
\end{itemize}

\subsubsection{Phase E: full ionization and thermal
closure}\label{phase-e-full-ionization-and-thermal-closure}

Port and validate:

\begin{itemize}
\tightlist
\item
  total ionization/recombination rates;
\item
  \passthrough{\lstinline!istruc!}/\passthrough{\lstinline!ioneqm!} ion
  fractions;
\item
  \passthrough{\lstinline!levwkelement!} source construction;
\item
  \passthrough{\lstinline!msolvelucy!} superlevel iteration;
\item
  heating/cooling and electron-fraction closure;
\item
  \passthrough{\lstinline!dsec!} convergence behavior.
\end{itemize}

\subsubsection{Phase F: radial transfer and output
products}\label{phase-f-radial-transfer-and-output-products}

Implement the zone/pass driver, continuum transfer, optical-depth
update, and output writers. Validate detail products first, then final
spectra:

\begin{lstlisting}
xoNN_detail/detal2/detal3/detal4
  -> xout_abund/lines/rrc/cont/spect
\end{lstlisting}

\subsubsection{Phase G: C++ backend}\label{phase-g-c-backend}

The Python implementation remains the readable reference and audit
layer. Performance-critical kernels should expose stable array-oriented
interfaces suitable for C++:

\begin{itemize}
\tightlist
\item
  cross-section mapping and
  \passthrough{\lstinline!phint53/phint53hunt!} integrals;
\item
  collision-rate grid evaluation;
\item
  sparse matrix assembly;
\item
  Lucy/sparse linear solves;
\item
  transfer and optical-depth loops;
\item
  thermal iteration.
\end{itemize}

The backend must consume and return the same typed Python state objects
and provenance IDs. Python tests and direct XSTAR probes remain the
correctness oracle.

\subsection{12. Conditional compact solve milestone (v0.3.201)}\label{conditional-compact-solve-milestone-v0.3.201}

For the six activated rows, partition the compact steady-state equations into selected and external blocks,
\begin{equation}
A_{SS}x_S + A_{SE}x_E = 0,
\qquad
A_{SS}x_S = -A_{SE}x_E.
\end{equation}
Here, $x_E$ is fixed to the captured XSTAR post-\texttt{msolvelucy} population vector. This isolates compact matrix assembly and source/sink closure without yet claiming an autonomous Python element solve. For O VII rows 79, 80, 241, 242, 244, and 293, the row-scaled selected matrix is full rank, its condition number is approximately 5.8545, and all six solved populations agree with XSTAR within 0.5\%. The maximum relative difference is approximately $3.0341\times10^{-3}$; the linear-system residual is below $4\times10^{-16}$.

This validates the selected compact topology and external source closure. The coefficients are still XSTAR-probed. The next implementation replaces them family by family with native Python rates, then expands the same partitioned solve to all 119 active rows before introducing a fully autonomous normalization and ionization closure.

\subsection{12.2. Native type-51 parity gate (v0.3.203)}

The first native family selected by the v0.3.202 readiness analysis is ATDB data type 51, the Burgess--Tully effective-collision-strength branch. XSTAR evaluates the spline at
\[
T_{\rm BT}=\max\!\left(T,\frac{2.8777\times10^6}{\lambda_{\AA}}\right),
\]
but retains the physical plasma temperature in the Maxwellian prefactor. For energy-ordered upper and lower levels,
\[
q_{u\rightarrow l}=\frac{8.626\times10^{-8}\,\Upsilon(T_{\rm BT})}{\sqrt{T/10^4\,\mathrm{K}}\,g_u},
\]
\[
q_{l\rightarrow u}=q_{u\rightarrow l}\,\frac{g_u}{g_l}\,
\exp\!\left[-\frac{\Delta E}{0.861707\,(T/10^4\,\mathrm{K})}\right],
\]
and
\[
\mathrm{ans1}=n_e q_{l\rightarrow u},\qquad
\mathrm{ans2}=n_e q_{u\rightarrow l}.
\]

The historical probe column named \texttt{xnx} was populated with the
\texttt{calc\_hmc\_ion.f90} variable \texttt{xee}; therefore the physical density is
\[
n_e=xpx\,xee.
\]
The v0.3.203 audit compares the native record rates and the four compact matrix insertions with the same XSTAR capture before enabling type-51 terms in the active compact solver.

\subsection{12.3. Native type-51 selected-system integration gate (v0.3.204)}

The v0.3.204 audit consumes the validated v0.3.199 record-term and population products together with the v0.3.203 native type-51 coefficients. For every compact term touching a selected row, it forms
\begin{equation}
A^{\mathrm{hybrid}}_{ij}=\begin{cases}
A^{\mathrm{native,51}}_{ij}, & \text{type 51},\\
A^{\mathrm{probe}}_{ij}, & \text{all other families}.
\end{cases}
\end{equation}
The replacement is applied wherever type 51 enters either the selected block $A_{SS}$ or the fixed-external closure $-A_{SE}x_E$. The all-probe and hybrid systems are reconstructed independently, their captured-population row residuals are compared, and both selected systems are solved.

The integration-ready flag requires exact one-use replacement of every selected-system type-51 term, coverage of every selected row, no missing population columns, a full-rank nonnegative solve, row-balance and population agreement within tolerance, and machine-level linear residuals. Type 50, type 71, type 53, and smaller families remain probe-backed, so complete native compact closure remains false and the production expanded-basis solver is unchanged.

\subsection{13. Validation policy}\label{validation-policy}

A feature becomes a production default only after passing all applicable
layers:

\begin{enumerate}
\def\labelenumi{\arabic{enumi}.}
\tightlist
\item
  \textbf{Database parity:} raw payload and level/ion identity are
  correct.
\item
  \textbf{Rate parity:} native \passthrough{\lstinline!ans1..ans6!}
  match the same XSTAR capture.
\item
  \textbf{Matrix parity:} row/column/sign and diagonal partners match.
\item
  \textbf{Basis parity:} compact index and alias roles match.
\item
  \textbf{Population parity:} pre/post solve vectors and normalization
  match.
\item
  \textbf{Emissivity parity:} lines/RRC/continuum match from the same
  populations.
\item
  \textbf{Cross-ion regression:} C V, O VII, Mg XI, and Ca XIX pass
  without ion-specific tuning.
\item
  \textbf{End-to-end parity:} radial/detail/final products agree within
  documented tolerances.
\end{enumerate}

\subsection{14. Practical source-to-Python
map}\label{practical-source-to-python-map}

\begin{longtable}[]{@{}
  >{\raggedright\arraybackslash}p{(\columnwidth - 2\tabcolsep) * \real{0.5000}}
  >{\raggedright\arraybackslash}p{(\columnwidth - 2\tabcolsep) * \real{0.5000}}@{}}
\toprule\noalign{}
\begin{minipage}[b]{\linewidth}\raggedright
XSTAR source
\end{minipage} & \begin{minipage}[b]{\linewidth}\raggedright
Python counterpart or target
\end{minipage} \\
\midrule\noalign{}
\endhead
\bottomrule\noalign{}
\endlastfoot
\passthrough{\lstinline!readtbl.f90!},
\passthrough{\lstinline!dread.f90!} &
\passthrough{\lstinline!inspect.py!},
\passthrough{\lstinline!hierarchy.ATDB!} \\
\passthrough{\lstinline!setptrs.f90!} &
\passthrough{\lstinline!ATDBIndexArrays!}, hierarchy cache,
process-specific joins \\
\passthrough{\lstinline!trnfrc.f90!},
\passthrough{\lstinline!bremsmap.f90!} &
\passthrough{\lstinline!RadiationField!}, live rate-grid probes, future
transfer backend \\
\passthrough{\lstinline!calc\_ion\_rates.f90!},
\passthrough{\lstinline!istruc.f90!} & current audits; future native
ionization closure \\
\passthrough{\lstinline!levwkelement.f90!} & basis/source probes; future
element-state constructor \\
\passthrough{\lstinline!calc\_hmc\_ion.f90!} & record-level parity,
compact matrix manifests \\
\passthrough{\lstinline!ucalc.f90!}, \passthrough{\lstinline!calt*.f90!}
& process-specific rate evaluators and provenance \\
\passthrough{\lstinline!phint53.f90!},
\passthrough{\lstinline!phint53hunt.f90!} & live type-53 audits; future
native PI/RR kernels \\
\passthrough{\lstinline!calc\_hmc\_element.f90!} & element-basis
probes/remap/scaffold/priority expansion \\
\passthrough{\lstinline!msolvelucy.f90!} &
\passthrough{\lstinline!xstar-lucy!} prototype and future
source-equivalent solver \\
\passthrough{\lstinline!calc\_emis*!},
\passthrough{\lstinline!linopac.f90!} & emissivity/export modules and
future opacity parity \\
\passthrough{\lstinline!stpcut.f90!} & future line/RRC depth
evolution \\
\passthrough{\lstinline!fstepr*!},
\passthrough{\lstinline!writespectra*!} & output readers now; future
writers \\
\end{longtable}

\subsection{15. References and source
provenance}\label{references-and-source-provenance}

Primary code reference: the XSTAR Fortran source distributed with the
reviewed package, especially \passthrough{\lstinline!readtbl.f90!},
\passthrough{\lstinline!setptrs.f90!},
\passthrough{\lstinline!xstarcalc.f90!},
\passthrough{\lstinline!calc\_hmc\_all.f90!},
\passthrough{\lstinline!calc\_hmc\_element.f90!},
\passthrough{\lstinline!calc\_hmc\_ion.f90!},
\passthrough{\lstinline!ucalc.f90!},
\passthrough{\lstinline!phint53.f90!},
\passthrough{\lstinline!calt68.f90!},
\passthrough{\lstinline!calt69.f90!},
\passthrough{\lstinline!calt71.f90!},
\passthrough{\lstinline!calt74.f90!},
\passthrough{\lstinline!calt77.f90!},
\passthrough{\lstinline!calt99.f90!},
\passthrough{\lstinline!msolvelucy.f90!},
\passthrough{\lstinline!trnfrc.f90!},
\passthrough{\lstinline!stpcut.f90!}, and the output routines.

Scientific references carried with the development review:

\begin{itemize}
\tightlist
\item
  Kallman et al.~(1996), \emph{ApJ}, 465, 994.
\item
  Kallman \& Bautista (2001), \emph{ApJS}, 133, 221.
\item
  Bautista \& Kallman (2001), \emph{ApJS}, 134, 139.
\item
  Kallman et al.~(2004), \emph{ApJS}, 155, 675.
\end{itemize}

The numerical findings in this document are from same-run XSTAR debug
probes and \passthrough{\lstinline!xstar-atomic!} audit products through
v0.3.201. They should be regenerated when the XSTAR source/database
version, benchmark inputs, or compact-basis selection changes.

\end{document}
